% GENERATED FILE. Edit canonical lessons, then run npm run generate:books.
% canonical-source-hash: fnv1a64:6957d5535ada803e
% canonical-lessons: ES-C702-mojarse, ES-C702-empapado, ES-C702-charco, ES-C702-barro, ES-C702-impermeable

\chapter{Mojarse}
\label{ch:es-mojarse}
\hlchaptermodality{702}

\begin{chapteropening}
\textbf{By the end of this chapter:} \emph{I can say I got caught in the rain: soaked, with puddles, mud and a raincoat.}

Complete the last lesson of chapter 702: I can say I got caught in the rain: soaked, with puddles, mud and a raincoat.
\end{chapteropening}

\section[mojarse]{mojarse — to get wet}

\label{lesson:ES-C702-mojarse}



\begin{quote}
Before the new one: say the Spanish for the drought, then the Spanish for to sweat.
\end{quote}

\subsection*{You'll want to know: mojarse}
\textbf{mojarse} — \textquotedblleft{}to get wet\textquotedblright{}.

\begin{quote}
Me mojé porque no llevaba paraguas.
\end{quote}

\begin{quote}
\emph{I got wet because I had no umbrella.}
\end{quote}

Without \emph{se}, \emph{mojar} is to wet something.

\begin{cousinweb}
\emph{mojarse} is built on \emph{mojar} \textquotedblleft{}to wet\textquotedblright{}, which is usually traced to Vulgar Latin \emph{molliāre} \textquotedblleft{}to soften\textquotedblright{}, from \emph{mollis} \textquotedblleft{}soft\textquotedblright{}: wet things go soft.
\end{cousinweb}

\begin{grammarlens}[title={What you've built}]
One more word for telling, explaining and describing at B1.
\end{grammarlens}

\subsection*{Your turn}
\begin{itemize}
\raggedright
  \item \emph{Say it:} \emph{mojarse}
  \item \emph{Say it:} \emph{mojarse}, once more
  \item \emph{Say it:} a sentence of your own with \emph{mojarse}
  \item \emph{Recall:} say the Spanish for the drought, then the Spanish for to sweat, then say \emph{mojarse} again
\end{itemize}

\begin{tcolorbox}[breakable,colback=teal!4,colframe=teal!35!black,title={Before you move on}]
What does mojarse mean? (\textquotedblleft{}To get wet\textquotedblright{}.) Say it once more.
\end{tcolorbox}
\section[empapado]{empapado — soaked, drenched}

\label{lesson:ES-C702-empapado}



\begin{quote}
Before the new one: say the Spanish for to sweat, then the Spanish for to get wet.
\end{quote}

\subsection*{You'll want to know: empapado}
\textbf{empapado} — \textquotedblleft{}soaked, drenched\textquotedblright{}.

\begin{quote}
Llegamos al hotel empapados.
\end{quote}

\begin{quote}
\emph{We got to the hotel soaked.}
\end{quote}

\emph{Empapado hasta los huesos} — soaked to the skin.

\begin{cousinweb}
\emph{empapado} is the past form of \emph{empapar} \textquotedblleft{}to soak\textquotedblright{}, which is built on \emph{papa} \textquotedblleft{}mush, pap\textquotedblright{}, from Latin \emph{pappa} \textquotedblleft{}food\textquotedblright{}, a word from baby talk (EtymDB).
\end{cousinweb}

\begin{grammarlens}[title={What you've built}]
One more word for telling, explaining and describing at B1.
\end{grammarlens}

\subsection*{Your turn}
\begin{itemize}
\raggedright
  \item \emph{Say it:} \emph{empapado}
  \item \emph{Say it:} \emph{empapado}, once more
  \item \emph{Say it:} a sentence of your own with \emph{empapado}
  \item \emph{Recall:} say the Spanish for to sweat, then the Spanish for to get wet, then say \emph{empapado} again
\end{itemize}

\begin{tcolorbox}[breakable,colback=teal!4,colframe=teal!35!black,title={Before you move on}]
What does empapado mean? (\textquotedblleft{}Soaked, drenched\textquotedblright{}.) Say it once more.
\end{tcolorbox}
\section[el charco]{el charco — the puddle}

\label{lesson:ES-C702-charco}



\begin{quote}
Before the new one: say the Spanish for to get wet, then the Spanish for soaked.
\end{quote}

\subsection*{You'll want to know: el charco}
\textbf{el charco} — \textquotedblleft{}the puddle\textquotedblright{}.

\begin{quote}
Los niños saltan en los charcos.
\end{quote}

\begin{quote}
\emph{The children jump in the puddles.}
\end{quote}

\emph{Cruzar el charco} — to cross the Atlantic.

\begin{cousinweb}
The origin of \emph{charco} is not settled; it may come from a pre-Roman word.
\end{cousinweb}

\begin{grammarlens}[title={What you've built}]
One more word for telling, explaining and describing at B1.
\end{grammarlens}

\subsection*{Your turn}
\begin{itemize}
\raggedright
  \item \emph{Say it:} \emph{el charco}
  \item \emph{Say it:} \emph{el charco}, once more
  \item \emph{Say it:} a sentence of your own with \emph{el charco}
  \item \emph{Recall:} say the Spanish for to get wet, then the Spanish for soaked, then say \emph{el charco} again
\end{itemize}

\begin{tcolorbox}[breakable,colback=teal!4,colframe=teal!35!black,title={Before you move on}]
What does el charco mean? (\textquotedblleft{}The puddle\textquotedblright{}.) Say it once more.
\end{tcolorbox}
\section[el barro]{el barro — the mud; clay}

\label{lesson:ES-C702-barro}



\begin{quote}
Before the new one: say the Spanish for soaked, then the Spanish for the puddle.
\end{quote}

\subsection*{You'll want to know: el barro}
\textbf{el barro} — \textquotedblleft{}the mud; clay\textquotedblright{}.

\begin{quote}
Tengo los zapatos llenos de barro.
\end{quote}

\begin{quote}
\emph{My shoes are covered in mud.}
\end{quote}

\emph{Una olla de barro} — a clay pot.

\begin{cousinweb}
\emph{barro} comes from Vulgar Latin \emph{barrum} \textquotedblleft{}mud, clay\textquotedblright{} (EtymDB); the deeper origin is not settled, and it may be pre-Roman.
\end{cousinweb}

\begin{grammarlens}[title={What you've built}]
One more word for telling, explaining and describing at B1.
\end{grammarlens}

\subsection*{Your turn}
\begin{itemize}
\raggedright
  \item \emph{Say it:} \emph{el barro}
  \item \emph{Say it:} \emph{el barro}, once more
  \item \emph{Say it:} a sentence of your own with \emph{el barro}
  \item \emph{Recall:} say the Spanish for soaked, then the Spanish for the puddle, then say \emph{el barro} again
\end{itemize}

\begin{tcolorbox}[breakable,colback=teal!4,colframe=teal!35!black,title={Before you move on}]
What does el barro mean? (\textquotedblleft{}The mud; clay\textquotedblright{}.) Say it once more.
\end{tcolorbox}
\section[el impermeable]{el impermeable — the raincoat; waterproof}

\label{lesson:ES-C702-impermeable}



\begin{quote}
Before the new one: say the Spanish for the puddle, then the Spanish for the mud.
\end{quote}

\subsection*{You'll want to know: el impermeable}
\textbf{el impermeable} — \textquotedblleft{}the raincoat; waterproof\textquotedblright{}.

\begin{quote}
Lleva el impermeable, va a llover.
\end{quote}

\begin{quote}
\emph{Take your raincoat, it is going to rain.}
\end{quote}

As an adjective: \emph{una chaqueta impermeable} — a waterproof jacket.

\begin{cousinweb}
\emph{impermeable} is a learned borrowing from Late Latin \emph{impermeābilis} \textquotedblleft{}that cannot be passed through\textquotedblright{} (EtymDB), built on \emph{meāre} \textquotedblleft{}to pass\textquotedblright{}.
\end{cousinweb}

\begin{grammarlens}[title={What you've built}]
One more word for telling, explaining and describing at B1.
\end{grammarlens}

\subsection*{Your turn}
\begin{itemize}
\raggedright
  \item \emph{Say it:} \emph{el impermeable}
  \item \emph{Say it:} \emph{el impermeable}, once more
  \item \emph{Say it:} a sentence of your own with \emph{el impermeable}
  \item \emph{Recall:} say the Spanish for the puddle, then the Spanish for the mud, then say \emph{el impermeable} again
\end{itemize}

\begin{tcolorbox}[breakable,colback=teal!4,colframe=teal!35!black,title={Before you move on}]
What does el impermeable mean? (\textquotedblleft{}The raincoat; waterproof\textquotedblright{}.) Say it once more.
\end{tcolorbox}
